(1) 设 \(f(x)=\tanh x\)，则
\(f'=\operatorname{sech}^2x\)，
\(f''=-2f\operatorname{sech}^2x\)，且 \(f^2-1=-\operatorname{sech}^2x\)，于是
\[
-\frac12f''+(f^2-1)f
=f\operatorname{sech}^2x-f\operatorname{sech}^2x=0.
\]
因此它是精确静止解。背景两端实幅度由 \(-1\) 变成 \(+1\)，相位差为 \(\pi\)。
相对于均匀凝聚的粒子数亏损
\(N_{\mathrm{def}}=\int(1-f^2)dx=\int\operatorname{sech}^2x\,dx=2\)。
能量两项各贡献 \(\frac12\operatorname{sech}^4x\)，因此
\(E=\int_{-\infty}^\infty\operatorname{sech}^4x\,dx=4/3\)。
无限背景本身能量已从题设能量泛函中扣除。

(2) 令 \(u,v\) 为实小扰动；保持一阶项，虚部和实部分别给
\[
u_t=-\frac12v_{xx},\qquad
-v_t=-\frac12u_{xx}+2u.
\]
用 \(u,v\propto e^{i(kx-\omega t)}\) 消元得
\[
\boxed{\omega_\pm^2=k^2+\frac{k^4}{4}.}
\]
按正频支 \(\omega_+(k)=|k|\sqrt{1+k^2/4}\)，
在 \(k>0\) 有
\[
v_{\rm ph}=\frac{\omega_+}k=\sqrt{1+k^2/4},\qquad
v_g=\frac{d\omega_+}{dk}=\frac{1+k^2/2}{\sqrt{1+k^2/4}}.
\]
故小波数声速 \(c_s=1\)。当 \(k<0\) 时正频群速为负，不能直接将只在正波数推导的相速符号用于负波数。

(3) 令 \(a=\sqrt{1-V^2}\)、\(y=a(x-Vt)\)。给定函数
\(\psi=iV+a\tanh y\) 满足
\[
\psi_t=-a^2V\operatorname{sech}^2y,\qquad
\psi_{xx}=-2a^3\tanh y\operatorname{sech}^2y,
\]
并有 \(|\psi|^2=1-a^2\operatorname{sech}^2y\)。于是
\(-\psi_{xx}/2+(|\psi|^2-1)\psi
=-iVa^2\operatorname{sech}^2y=i\psi_t\)，确为精确解。
中心最小密度 \(V^2\)，密度亏损
\[
N_{\rm def}(V)=a^2\int_{\mathbb R}\operatorname{sech}^2[a(x-Vt)]dx=2a.
\]
两端取值 \(-a+iV\)、\(a+iV\)，对正 \(V\) 取沿短弧连续的相位差绝对值为 \(2\arccos V\)；负速度时选对应相位支。\(V=0\) 回到黑孤子，\(|V|\to1\) 时趋于均匀密度、亏损趋零。

(4) 对灰孤子，
\(|\psi_x|^2=a^4\operatorname{sech}^4y\) 与
\((|\psi|^2-1)^2=a^4\operatorname{sech}^4y\)，故
\[
E(V)=a^4\int_{\mathbb R}\operatorname{sech}^4[a(x-Vt)]dx
=\frac43 a^3=\frac43(1-V^2)^{3/2}.
\]
当 \(0<V<1\)，
\(dE/dV=-4V\sqrt{1-V^2}<0\)。
黑/灰孤子是非线性与色散自洽平衡的精确移动结构，因此在无限、均匀、无耗散的一维模型中以恒速传播并保持局域形状；外势、有限温度、横向激发和不可积损耗都会改变该结论。
